\section{Controlled Ablations}
\paragraph{Past scene information.}
In the exposed own-B experiment, comparing motion-and-geometry with
quality-and-geometry holds the requested prediction support fixed.
Table~\ref{tab:context} gives the reported free-switch policies; their
realized intervention rates differ, so the table is not an equal-count
effect. The registered direct motion-minus-quality conservative
contrast is \ContextContrast. It does not establish a scene-input
benefit. The motion policy improves hard-slice available-label error
while worsening the easy slice and still fails the complete risk rule.
In a selected worst-easy case with all twelve labels present, retaining
CV gives \FailureBaseline{} pixels ADE while the chosen neural
forecast gives \FailureCandidate{} pixels ADE. This is a tail case,
not an estimate of typical behavior.

\begin{table*}[t]
\centering\small
\caption{Past-image cost inputs on exposed own-B development.
Accuracy columns report available-label gain over CV in percent.
The 72 views repeat 24 recordings across three risk-head seeds.}
\label{tab:context}
\begin{tabular}{@{}lrrrr@{}}
\toprule
Input & All & Hard & Easy & Original risk passes / 72 \\
\midrule
\ContextRows
\bottomrule
\end{tabular}
\end{table*}

\paragraph{Interaction ablation.}
With matched forecasters, query counts and the same observed-neighbor
graph, the relative-trajectory ADE reduction of joint versus additive
is \PairObserved{} pixels on observed pairs. The full-grid conservative
lower contrast is \PairLower{} pixels; the upper contrast is
\PairUpper{} pixels. Thus the observed pair comparison does not
establish an improvement with missing pair outcomes, nor does it
measure physical collisions. The conservative trajectory utility
contrast in Table~\ref{tab:matched}'s setting likewise does not
support the added cross term.

\paragraph{Risk after selection.}
The local screen accepts options according to predictions fitted on
candidate actions. Refitting to the actions the search actually selects
tests whether this mismatch can be repaired. With 3,280 matched
switches, learned probability ranking versus the matched margin
control gives \PostRiskContrast; both pass one of 102 views.
A separate loss directly fits additive harm and reference masses,
then applies a past-only recording ledger. At 280 switches matched
within recording and cardinality, whole-recording training gives
\PostMassContrast{} relative to the pointwise margin rank. The
recording model passes three views and the matched control four.
These fixed-count ranks are retrospective controls, not deployments.
Neither of the nested post-search model families has a threshold
that passes every required validation view in any of its 42 folds.

\paragraph{Tail errors and support.}
The original gated joint strategy uses \GateSwitches{} switches in
the native Moving A experiment; the CV and gated joint worst
recording-level p95 available-label ADE are both \TailCV{} pixels.
Unconditionally substituting OLS4 yields \TailOLS{} pixels in this
tail summary. OLS4 is a comparison forecast, not a policy certified
against CV. CV abstains everywhere, so its selected easy risk is
undefined, not an observed pass. The forecast role and the risk
denominator must remain explicit when interpreting either comparator.
